\chapter{Requirements Analysis}
\label{chap:requirements}

This chapter specifies the functional and non-functional requirements of the system, the hardware used for development and deployment, and the software stack.

\section{Functional Requirements}

\begin{table}[H]
\centering
\caption{Functional requirements}
\label{tab:functional_requirements}
\begin{tabularx}{0.95\textwidth}{>{\bfseries}l X}
\toprule
ID & Requirement \\
\midrule
FR-01 & Generate reference skyline profiles from a DEM for a given study area. \\
FR-02 & Accept one or more overlapping photographs as input. \\
FR-03 & Segment sky from terrain in each photograph. \\
FR-04 & Extract a 1-D elevation-angle profile from the segmented mask. \\
FR-05 & Match the profile against a reference database to estimate camera location. \\
FR-06 & Return a confidence status: \textit{OK}, \textit{LOW\_CONFIDENCE}, \textit{NO\_MATCH}, or \textit{INVALID\_INPUT}. \\
FR-07 & Capture device sensor data (pitch, heading) alongside each photo. \\
FR-08 & Fuse multiple overlapping photographs into a single wide-FOV profile. \\
\bottomrule
\end{tabularx}
\end{table}

\section{Non-Functional Requirements}

\begin{table}[H]
\centering
\caption{Non-functional requirements}
\label{tab:nonfunctional_requirements}
\begin{tabularx}{0.95\textwidth}{>{\bfseries}l X}
\toprule
ID & Requirement \\
\midrule
NFR-01 & Work offline, no internet. \\
NFR-02 & Lightweight segmentation model for mobile. \\
NFR-03 & Stream the reference database in small batches (peak \Gls{ram} $<$ 150\,MB). \\
NFR-04 & Database regenerable for different study areas by swapping the DEM. \\
NFR-05 & Modular pipeline, each component updatable independently. \\
NFR-06 & Segmentation, extraction, and matching on-device via ONNX. \\
NFR-07 & Deterministic output for the same input and database. \\
\bottomrule
\end{tabularx}
\end{table}

\section{Development Hardware}

\begin{table}[H]
\centering
\caption{Development hardware}
\label{tab:dev_hardware}
\begin{tabular}{ll}
\toprule
\textbf{Component} & \textbf{Specification} \\
\midrule
Processor & Multi-core x86-64 CPU \\
Memory & 16\,GB RAM \\
Storage & 50\,GB free (DEMs, database, training data) \\
GPU & NVIDIA T4 (Google Colab free tier) \\
OS & Linux (Ubuntu 22.04) \\
\bottomrule
\end{tabular}
\end{table}

\section{Target Hardware (Mobile)}

\begin{table}[H]
\centering
\caption{Target mobile hardware}
\label{tab:target_hardware}
\begin{tabular}{ll}
\toprule
\textbf{Component} & \textbf{Minimum} \\
\midrule
OS & Android 10 (API 29) \\
Architecture & arm64-v8a \\
RAM & 3\,GB \\
Storage & 100\,MB free (ONNX model + DB subset) \\
Camera & Rear, autofocus \\
Sensors & Accelerometer, magnetometer \\
\bottomrule
\end{tabular}
\end{table}

\section{Software Stack}

\begin{table}[H]
\centering
\caption{Software stack}
\label{tab:software_stack}
\begin{tabularx}{0.95\textwidth}{l X}
\toprule
\textbf{Tool} & \textbf{Used for} \\
\midrule
Python 3.10 & Pipeline code, evaluation scripts \\
PyTorch & Training the U-Net \\
segmentation-model-pytorch & U-Net wrapper with pretrained MobileNetV3 \\
Albumentations & Training augmentation (flips, colour jitter) \\
OpenCV & Image I/O, CLAHE, Canny \\
NumPy \& SciPy & Arrays, FFT cross-correlation, Gaussian filtering \\
Rasterio & Reading GeoTIFF DEMs \\
HORAYZON & Ray-tracing horizon profiles from DEM mesh \\
PyRender & Rendering synthetic scenes (headless via Embedded Graphics Library) \\
trimesh & DEM vertices to triangle mesh \\
FastDTW & Fine alignment between query and reference \\
pyarrow \& pandas & Parquet database I/O \\
geopy & Geodesic distances for error metrics \\
Kivy & Mobile UI framework \\
Plyer & Accelerometer/magnetometer access \\
ONNX Runtime & On-device segmentation inference \\
Git & Version control \\
\LaTeX & Report \\
\bottomrule
\end{tabularx}
\end{table}
